% Eysenck's idea: introverts start with higher cortical arousal, extroverts lower, so extroverts seek
% extra stimulation to reach a comfortable level. Schematic, not data. The extrovert's top-up shows on click 3.
\begin{tikzpicture}[lbl/.style={font=\small, text=plum, align=center}]
  % comfort line
  \draw[plum!70, dashed, line width=0.9pt] (-2.1,2.9) -- (2.3,2.9);
  \node[anchor=east, font=\footnotesize, text=plum, align=right] at (-2.15,2.9) {հարմարավետ\\ մակարդակ};
  % introvert tube
  \fill[lavender, rounded corners=2mm] (-1.9,0) rectangle (-0.9,2.65);
  \draw[plum, line width=0.9pt, rounded corners=2mm] (-1.9,0) rectangle (-0.9,3.6);
  \node[lbl, font=\blackfont\small] at (-1.4,-0.35) {Ինտրովերտ};
  \node[lbl, font=\footnotesize] at (-1.4,-0.85) {արդեն գրեթե\\ բավական է};
  % extrovert tube
  \fill[peach, rounded corners=2mm] (0.9,0) rectangle (1.9,1.2);
  \fill[amOrange!70, visible on=<3->] (0.9,1.2) rectangle (1.9,2.9);
  \draw[plum, line width=0.9pt, rounded corners=2mm] (0.9,0) rectangle (1.9,3.6);
  \node[lbl, font=\blackfont\small] at (1.4,-0.35) {Էքստրավերտ};
  \node[lbl, font=\footnotesize] at (1.4,-0.85) {պետք է\\ լրացնել};
  \node[visible on=<3->] at (2.7,2.0) {\bigemoji[0.75cm]{party}};
  \node[visible on=<3->] at (2.7,1.1) {\bigemoji[0.75cm]{microphone}};
\end{tikzpicture}
